\begin{frame}{더티 데이터 (끝): 이상치 --- 도메인 지식의 영역}
  \begin{itemize}
    \item 3. \kb{이상치 (outlier)}: 범위는 합리적이지만 분포와
      \textbf{극단적으로 어긋난} 값.
    \item 음수 가격, 250세 나이 = 센티넬 $\to$ \textbf{1번(결측치)}으로.
      ``합법적이지만 극단적''인 값(수백억짜리 단독주택 한 채)은
      \textbf{3번(이상치)}으로 처리할지, 드물지만 \textbf{중요한 사례}
     라면 그대로 둘지는 ---
  \end{itemize}
  \vskip0.4em
  \begin{block}{판단의 성격}
    \kq{도메인 지식으로만} 판단할 수 있는 문제.
    ``코드 짜는 것''이 아니라 ``\textbf{문제를 이해하는 것}''의
    영역에 속한다는 점을 기억할 것.
  \end{block}
  \vskip0.4em
  \begin{itemize}
    \item 세 모양 모두 \textbf{정제 단계}에서 다루고, 발견은
      \textbf{EDA 단계}에서 한다 --- 다음 4가지 질문으로.
  \end{itemize}
\end{frame}
